\section{Introduction}\label{sec:intro}

\subsection{The problem}
Chebyshev observed in 1853 that primes $\equiv3\pmod4$ seem to outnumber primes $\equiv1\pmod4$. Rubinstein and Sarnak \cite{RS} made this precise: under the generalized Riemann hypothesis (GRH) for $L(s,\chi_{-4})$ and the linear independence hypothesis (LI: the positive ordinates of the zeros are linearly independent over $\QQ$), the set $\{x:\pi(x;4,3)>\pi(x;4,1)\}$ has logarithmic density $0.9959\ldots$. GRH alone already gives a limiting logarithmic distribution $\nu$ of the normalised race $(\log x/\sqrt x)(\pi(x;4,3)-\pi(x;4,1))$ \cite{RS}, which is the push-forward of Haar measure on the closure of a one-parameter subgroup of an infinite-dimensional torus (\cite{ANS,Devin,Bailleul}; \cite[Theorem~1.3]{GO}). The density exists if $\nu(\{0\})=0$, and LI is used to exclude an atom at $0$. Rubinstein and Sarnak note (\cite[p.~174]{RS}) that ``assuming only GRH we don't know that $\delta(P_{q;a_1,\dots,a_r})$ exists''. As far as we know this is still open. Several weakenings of LI suffice (Devin \cite[Theorems~2.2, 5.4]{Devin}, Martin--Ng \cite{MN}, Bailleul \cite{Bailleul}); none of them is known to hold. Item~1 of the problem list of Hamieh, Kadiri, Martin and Ng \cite{HKMN} (posed by Martin and Silberman) asks to prove that $\zeta(s)$ has at least one irrational ordinate, and that the positive ordinates of $\zeta(s)$ are not all rational multiples of one another; the analogous statements for $L(s,\chi)$ are equally open. Problem~13 of the same list asks to prove unconditionally that ties $\pi(x;q,a)=\pi(x;q,b)$ have density $0$; our results bear on this only conditionally on GRH (Corollary~\ref{cor:ties}). Under GRH and weakened forms of LI, density $0$ for ties follows from the continuity results of Devin \cite[Theorem~2.2]{Devin}, \cite[Corollary~2]{Devin2}.

\textbf{We do not prove that the logarithmic density of $\{x:\pi(x;4,3)>\pi(x;4,1)\}$ exists under GRH alone.} This paper collects what we could prove in that direction, and why the methods available to us cannot finish the job.

\subsection{Setting}
Let $\chi$ be a real primitive character modulo $q$, $\chi^*=-\chi$, and for real $m>-\frac12$ put $M=m+\frac12$ and
\[
D_m(x)=\sum_{p\le x}\chi^*(p)p^m,\qquad E_M(x)=2M\,x^{-M}(\log x)\,D_m(x).
\]
For $q=4$, $m=0$: $D_0(x)=\pi(x;4,3)-\pi(x;4,1)$ and $E_{1/2}(x)=(\log x/\sqrt x)D_0(x)$; for $\chi=\chi_{-3}$, $D_0(x)=\pi(x;3,2)-\pi(x;3,1)$. The weight $p^m$ interpolates between the classical race ($m=0$) and the race weighted by $p^{-1/2}$ ($m\to-\frac12$), for which Hayani \cite[Corollary~1.2]{Hayani} proved, under GRH alone, that the lead set has natural density $1$ ($q=4$); Sheth \cite{Sheth} has closely related results under GRH, and the $p^{-1/2}$-weighted framework goes back to Aoki and Koyama \cite{AK}. GRH$(\chi)$ means that all nontrivial zeros of $L(s,\chi)$ have real part $\frac12$ (no hypothesis on any other $L$-function, none on $\zeta$); NCZ$(\chi)$ means $L(\frac12,\chi)\ne0$, which is a theorem for $q=3,4$. Under GRH and NCZ, $E_M(e^y)$ has a limiting distribution $\nu_M$ in $y$ (Section~\ref{sec:setting}). Results conditional on GRH say so in their statements. A Lean proof of GRH for Dirichlet $L$-functions was announced on 2026-11-20; it is unrefereed and is not used anywhere in this paper. It would in any case not touch the question studied here, which concerns linear relations among the zeros.

\subsection{Results}
All statements below assume GRH$(\chi)$ and NCZ$(\chi)$ unless stated otherwise, and none assumes LI.

\emph{Certified bounds} (Theorem~\ref{thm:cert}). Under GRH$(\chi_{-4})$ alone, the lower logarithmic density of $\{x:\pi(x;4,3)>\pi(x;4,1)\}$ is $>0.910924861$; under GRH$(\chi_{-3})$ alone, that of $\{x:\pi(x;3,2)>\pi(x;3,1)\}$ is $>0.934761398$. The previous certified values without LI were $0.884276434$ and $0.911240824$ \cite{GO}. So $\nu_{1/2}(\{0\})<0.0891$, resp.\ $<0.0653$.

\emph{Time reversal} (Theorem~\ref{thm:reversal}). On the closure torus $\cA$, with $F_M$ the function whose law is $\nu_M$, one has $F_M(\theta)+F_M(\theta^{-1})=2+R(\theta)$ almost everywhere, with $|R|\le\sigma_M^2=\sum_{\gamma>0}\frac{8M^2}{M^2+\gamma^2}$ (sum over the positive ordinates with multiplicity; a closed form). So $\nu_M$ is symmetric about its mean $1$ up to $\sigma_M^2$ in the $L^\infty$-Wasserstein distance, without LI, and $\nu_M((-\infty,0])\le\frac12$ whenever $\sigma_M^2<2$.

\emph{Atom $=$ jump} (Theorem~\ref{thm:jump}). For $M$ in the interior of an explicit interval containing $\frac12$, $\nu_M(\{0\})$ equals the jump at $m$ of the nonincreasing function $m\mapsto\dlow\{x:D_m(x)>0\}$. Hence the logarithmic density of $\{D_m>0\}$ exists for all but countably many $m\in(-\frac12,1.3239]$ ($q=4$), and the classical race has a density as soon as the bias is continuous in the weight at $m=0$.

\emph{Sufficient conditions} (Section~\ref{sec:sufficient}). The density exists for every $m$ under GRH together with any one of the following hypotheses on the ordinates, each implied by LI: TS (the ordinates above some height do not span all ordinates over $\QQ$); FD (finite-level directions); TS$_\ZZ$ (the ordinates above height $T$ generate subgroups of unbounded index); V (some ordinate never carries half of the $\ell^1$-mass of an integer relation); GLI (graded linear independence: relations reaching height $t$ have length $\ge\psi(t)$ with $\psi(t)\log t/t\to\infty$); or a prime-side band-decoupling hypothesis on mixed moments of windowed prime sums (implied by LI together with simplicity of the zeros). As conditions on sets of frequencies, TS and V are strictly weaker than the hypothesis of Devin's Theorem~2.2, and FD is strictly weaker than each hypothesis of this kind that we found in the literature. Conversely, an atom forces strong arithmetic structure: every ordinate has a multiple $n\gamma$, $n\le1/\nu_M(\{c\})$, that is an integer combination of ordinates above any given height, and integer relations of length $\le Ct/\log t$ reaching height $t$ occur infinitely often.

\emph{Reduction} (Lemma~\ref{lem:reduction}, Conjecture~\ref{conj:U}). The density at $m=0$ follows from a uniform small-ball bound for finite trigonometric polynomials with positive integer multiplicities obeying the zero-counting law, on arbitrary subtori; in this form the unknown group generated by the zeros disappears. We prove this bound on the subclass of models with a low Newton vertex (Corollary~\ref{cor:Ucase}).

\emph{Barriers} (Section~\ref{sec:barriers}). There are models with all the soft structure used above (a compact connected hull with dense flow, positive weights, the identities between weights, time reversal, $\sigma^2<2$ at the classical weight) whose law has an atom at $0$ (Proposition~\ref{prop:model73}, certified). Integer multiplicities forbid flat arcs (Theorem~\ref{thm:arcs}) and forbid atoms in structured classes (Theorem~\ref{thm:nested} and Section~\ref{sec:integrality}), but not in general: there are positive integer multiplicities $m_n\ll(\log n)^u$, for any $u>0$, whose model vanishes on a set of positive measure, and at the classical weight such models can have $\sigma^2<\frac12$ and obey the local zero-counting law, with a constant of size at least about $10^{14}$ (larger for the models with $\sigma^2<\frac12$) (Theorem~\ref{thm:noise}, Corollary~\ref{cor:noise}). The fixed-constant form (Conjecture~\ref{conj:U}) and the case of bounded multiplicities are untouched by this.

\subsection{Relation to the literature}
That a race against a \emph{level} has a density for all but countably many levels is classical; Devin \cite{Devin} attributes a statement of this type under RH alone to Kaczorowski and Ramar\'e \cite[Theorem~3]{KR} (we have not read it). What is specific to Theorem~\ref{thm:jump} is that the level is fixed at $0$ and the weight varies; it extends \cite[Theorem~7.1]{GO}, which gave only that ties are immaterial outside a countable set of weights. The mechanism of Theorems~\ref{thm:TS}, \ref{thm:FD} (move along a one-parameter subgroup of the hull on which the function is a non-constant trigonometric polynomial) is Devin's \cite{Devin}; what is added is that conditioning on the complementary coordinates replaces independence. The nearest prior result to Theorem~\ref{thm:reversal} that we found is Devin's exact symmetry \cite[Theorem~2.3]{Devin} under a parity hypothesis on the integer relations among the ordinates (also unproved). Chebyshev-type lower bounds without LI, with the multiplicity-squared variance, are Devin \cite[Corollary~5.8]{Devin} and Fiorilli \cite[Lemma~2.7]{Fiorilli}. In the papers we searched (listed in \file{density/LITERATURE.md}) we did not find the time-reversal identity, an explicit GRH-only lower bound for the bias of the $q=4$ race other than the author's \cite{GO}, or hypotheses of the TS, FD, V type; this is a statement about those searches only. No search was made for finite-index conditions or graded independence conditions, and no novelty is claimed for TS$_\ZZ$ and GLI. Lemma~\ref{lem:JN} is classical.

\subsection{Status of the results}
Every result stated as a theorem, proposition, lemma or corollary passed at least one adversarial read by an independent AI referee instructed to find errors (Section~\ref{sec:review}); the number of reads is given in brackets in each statement, e.g.\ \rd{2 reads}. Statements labelled \emph{Remark} or \emph{Heuristic} are not claimed as proved, except where a remark explicitly derives a consequence of a stated result. Many proofs are given in full; many others (in particular in Sections~\ref{sec:sufficient}--\ref{sec:barriers}) are outlined or only referenced, with complete proofs in the working notes cited at each place. File names refer to the directory \repo{} of \cite{Repo} unless they start with \texttt{05-}.

Table~\ref{tab:summary} lists the main results.

\begin{table}[ht]\small
\caption{Main results. ``GRH'' means GRH$(\chi)$ and NCZ$(\chi)$ for the character concerned (NCZ is a theorem for $q=3,4$).}\label{tab:summary}
\begin{tabular}{p{0.37\linewidth}p{0.2\linewidth}p{0.12\linewidth}p{0.17\linewidth}}\toprule
Statement&Hypotheses&Where&Reads\\\midrule
certified lower densities $0.910924861$ ($q=4$), $0.934761398$ ($q=3$)&GRH&Thm~\ref{thm:cert}&2, certified\\
time reversal; $\nu_M((-\infty,0])\le\frac12$ if $\sigma_M^2<2$&GRH&Thm~\ref{thm:reversal}&2\\
atom at $0$ $=$ jump in the weight; density for all but countably many weights&GRH&Thm~\ref{thm:jump}&2\\
density exists under TS, FD, tame directions&GRH $+$ hypothesis&\S\ref{sec:TS}&1\\
density exists under TS$_\ZZ$&GRH $+$ TS$_\ZZ$&Thm~\ref{thm:TSZ}&2\\
density exists under V; low vertices&GRH $+$ V / LV&Thm~\ref{thm:V}&1\\
density exists under GLI; what an atom forces&GRH $+$ GLI&Thm~\ref{thm:GLI}&2\\
prime-side duality; band decoupling&GRH $+$ prime-side&Thms~\ref{thm:P1}, \ref{thm:P7}&2\\
reduction to finite integer models; U on a subclass&GRH / none&\S\ref{sec:reduction}&1\\
positive-weight models with an atom, $M=\frac12$, $\sigma^2<0.55$&none&Prop~\ref{prop:model73}&1, certified\\
integrality: parity on nested combs, Toeplitz classes, arcs&none&\S\ref{sec:integrality}&1--2\\
integer models with an atom, $m_n\ll(\log n)^u$, counting constant $\gtrsim10^{14}$&none&Thm~\ref{thm:noise}&1\\\bottomrule
\end{tabular}
\end{table}
